\documentclass[11pt]{article}
\usepackage[margin=1in]{geometry}
\usepackage{amsmath,amssymb,mathtools}
\usepackage{amsthm}
\usepackage{booktabs,array,longtable}
\usepackage[hidelinks]{hyperref}
\usepackage{xurl}
\usepackage{xcolor}
\usepackage{fontspec}
\usepackage{microtype}

% XeLaTeX preserves the UTF-8 spelling of Cubical Agda identifiers.
% In particular, \Agda keeps π₁, ≃∙, ≡, ↑, and → as literal code characters.
\setmonofont{Menlo}
% Arial Unicode MS supplies the mathematical and subscript glyphs absent from
% Menlo while retaining the source's literal UTF-8 identifiers.
\newfontfamily\AgdaFont{Arial Unicode MS}[Scale=MatchLowercase]
\renewcommand\UrlFont{\AgdaFont\upshape}
\newcommand{\Agda}[1]{{\normalfont\nolinkurl{#1}}}
\newcommand{\AgdaOptions}{{\AgdaFont\char123\#- OPTIONS --cubical --guardedness \#-\char125}}
\newcommand{\BG}{\mathsf{BG}_{\mathrm{pres}}}
\newcommand{\FP}{F_P}
\newcommand{\GP}{G_P}
\newcommand{\piOne}{\pi_1}
\setlength{\emergencystretch}{3em}
\setlength{\parskip}{0.35em}
\setlength{\parindent}{1.5em}
\newcolumntype{L}[1]{>{\raggedright\arraybackslash}p{#1}}
\newtheorem{definition}{Definition}[section]
\newtheorem{lemma}[definition]{Lemma}
\newtheorem{proposition}[definition]{Proposition}
\newtheorem{theorem}[definition]{Theorem}
\newtheorem{corollary}[definition]{Corollary}
\newtheorem{remark}[definition]{Remark}
\newcommand{\formal}[1]{\par\noindent\footnotesize\normalfont\textit{Formal source.} \Agda{#1}.\par\normalsize\smallskip}
\title{Group Presentations as Homotopy Pushouts\\
\large A Cubical Agda Formalization}
\author{
Hongwei (Alex) Wang\\[3pt]
\small New York University Abu Dhabi\\
\small \texttt{hw4375@nyu.edu}
}
\date{11 November 2026}

\begin{document}
\maketitle

\begin{abstract}
For a presentation by a type of generators and a type of relations, we form a
1-truncated homotopy pushout obtained by attaching one circle for each
relation.  We construct the usual algebraic quotient of the free group by the
normal closure of the relation differences and a canonical homomorphism from
this group to the fundamental group of the resulting type.  The principal
formal result is that this homomorphism is surjective.  The article gives a
mathematical account of the checked Cubical Agda proofs; it does not assert
injectivity, a presentation theorem for the fundamental group, or a
classifying-space equivalence.
\end{abstract}


\section{Introduction}

Let $G$ and $R$ be types.  A presentation with generators $G$ and relations
$R$ may be realized geometrically by a graph with one loop for each generator,
followed by the attachment of a circle for each relation.  In the present
construction the graph is expressed as a suspension, and the attachment is a
homotopy pushout.  The subsequent 1-truncation retains the range in which a
fundamental group may be formed.  The construction is carried out in Cubical
Agda, using its higher inductive types for suspensions, circles, pushouts, and
truncations~\cite{cubical,hottbook}.

The input is universe-polymorphic: $G:\mathsf{Type}_{\ell}$,
$R:\mathsf{Type}_{\ell'}$, and a presentation $P$.  No sethood,
decidability, or discreteness assumption is made on $G$ or $R$.  The source
annotations below identify the declaration whose type and proof term establish
the assertion in question.  Environment and API details are deferred to
Appendix~A.

\section{Presentations and loop words}

\formal{GroupPresentation.Presentation; GroupPresentation.Signed; GroupPresentation.Word; GroupPresentation.letterLoop; GroupPresentation.wordLoop; GroupPresentation.relationLoop₀; GroupPresentation.relationLoop₁; GroupPresentation.relationLoop}
\begin{definition}
A \emph{syntactic presentation} $P$ on $G$ and $R$ consists of two functions
$p_0,p_1:R\to\mathsf{Word}(G)$, where a word is a finite list of signed
generators.  Once a based space carrying loops $\ell_g$ is fixed, a positive
letter is evaluated by $\ell_g$, a negative letter by $\ell_g^{-1}$, and a
word by concatenation.  Write $L_i(r)$ for the loop represented by $p_i(r)$
and $L(r)=L_0(r)\cdot L_1(r)^{-1}$ for its relation-difference loop.
\end{definition}

The words are evaluated in a raw identity type, before set truncation.  Thus
the circle attachment uses the actual path $L(r)$ and requires no normal-form
theorem for an untruncated loop space.

\section{The generator space}

\formal{GroupPresentation.GeneratorInput; GroupPresentation.GeneratorSpace; GroupPresentation.generatorLoop}
\begin{definition}
Let $e$ denote the point of the auxiliary summand $\mathbf{1}$.  The
\emph{generator space} is $X_G=\mathrm{Susp}(\mathbf{1}\sqcup G)$, based at
its north pole.  For $g:G$, put $\ell_g=\mathrm{merid}(g)\cdot
\mathrm{merid}(e)^{-1}$.  Thus each generator determines a north-based loop.
\end{definition}

\formal{GroupPresentation.Connectedness.generatorSpaceMerePathToNorth}
\begin{lemma}
Every point of $X_G$ is merely equal to its north pole.
\end{lemma}
\begin{proof}
The proof is suspension induction.  At north it is represented by reflexivity,
and at south by the reference meridian.  Along a meridian the target is a
proposition, namely a propositional truncation of an identity type; the two
endpoint representatives therefore extend uniquely.
\end{proof}

\formal{GroupPresentation.GeneratorBouquet.generatorSpace≃∙Bouquet; GroupPresentation.GeneratorBouquet.generatorLoop-toBouquet; GroupPresentation.GeneratorBouquet.liftedGeneratorSpace≃∙Bouquet}
\begin{proposition}
There is a pointed equivalence $X_G\simeq_*\mathrm{Bouquet}(G)$ sending
$\ell_g$ to the $g$-th bouquet loop; the same assertion holds after the
universe lift used below.
\end{proposition}
\begin{proof}
The forward suspension eliminator sends both poles to the bouquet base point,
the $e$-meridian to reflexivity, and the $g$-meridian to the $g$-loop.  The
inverse sends the bouquet base point to north and its $g$-loop to $\ell_g$.
Their homotopies are respectively suspension and bouquet induction.  The
loop calculation follows by functoriality of concatenation and the unit law
for the image of the reference meridian.  The lifted proof is the same
construction with $G$ replaced by its lift.
\end{proof}

\section{Relation circles and the presentation type}

\formal{GroupPresentation.attachingMap; GroupPresentation.relationProjection; GroupPresentation.PresentationPushout; GroupPresentation.BG_pres; GroupPresentation.BG-basepoint}
\begin{definition}
For $r:R$, circle recursion sends the circle base point to north and its
generating loop to $L(r)$.  Let $a:R\times S^1\to X_G$ be the resulting map
and let $q:R\times S^1\to R$ be projection.  Define
\[
Y_P=\operatorname{Pushout}(a,q), \qquad B_P=\lVert Y_P\rVert_1.
\]
The formal notation is \Agda{PresentationPushout P} and \Agda{BG_pres P};
the latter is implemented as h-level truncation $3$, the library's index for
a 1-type.
\end{definition}

\formal{GroupPresentation.Relations.relationCircleNaturality; GroupPresentation.Relations.conjugacyToRefl; GroupPresentation.Relations.inlRelationLoop-null; GroupPresentation.Relations.inlRelationSidesAgree}
\begin{proposition}
For every $r:R$, the image of $L(r)$ in $Y_P$ is null.  Consequently the
images of $L_0(r)$ and $L_1(r)$ in $Y_P$ are equal.
\end{proposition}
\begin{proof}
Write $u$ for the pushout glue path at $(r,\mathsf{base})$.  Naturality of
the family of glue paths around the circle gives $u=\mathsf{inl}(L(r))\cdot
u$.  Cancelling $u$ on the right proves $\mathsf{inl}(L(r))=\mathsf{refl}$;
the cited term performs this cancellation explicitly by the right unit,
associativity, and inverse-cancellation laws.  Since $L(r)=L_0(r)\cdot
L_1(r)^{-1}$, a second cancellation calculation gives
$\mathsf{inl}(L_0(r))=\mathsf{inl}(L_1(r))$.
\end{proof}

\formal{GroupPresentation.TruncatedRelations.generatorSpaceToBG; GroupPresentation.TruncatedRelations.BG-relationLoop-null; GroupPresentation.TruncatedRelations.BG-relationSidesAgree}
\begin{corollary}
The canonical map $X_G\to B_P$ sends every relation loop to reflexivity and
sends the two sides of each relation to equal loops.
\end{corollary}
\begin{proof}
The map is the composite of the left pushout inclusion with the truncation
constructor.  Applying that constructor to the two preceding path equalities
gives the asserted equalities in $B_P$.
\end{proof}

\formal{GroupPresentation.Connectedness.presentationPushoutMerePathToBase; GroupPresentation.Connectedness.BG-pres-isConnected}
\begin{proposition}
$B_P$ is connected.
\end{proposition}
\begin{proof}
Pushout induction reduces a point of $Y_P$ either to a point of $X_G$, which
is merely joined to north by the preceding lemma, or to a relation point,
which is joined to the base point by the glue path at the circle base point.
Truncation induction transports these merely inhabited path spaces to $B_P$.
The final term contracts the set truncation at the selected base point, which
is the library definition of \Agda{isConnected 2}.
\end{proof}

\section{The algebraic presented group}

\formal{GroupPresentation.PresentedGroup.PresentationFreeGroup; GroupPresentation.PresentedGroup.wordToFreeGroup; GroupPresentation.PresentedGroup.relationDifference; GroupPresentation.PresentedGroup.presentationNormalSubgroup; GroupPresentation.PresentedGroup.PresentedGroup}
\begin{definition}
Put $F_P=\mathsf{FreeGroup}(\mathrm{Lift}_{\ell'}G)$. A signed word is evaluated in $F_P$, and $d_r=[p_0(r)][p_1(r)]^{-1}$. Let $N_P$ be the normal subgroup generated by the lifted family $(d_r)_{r:R}$, and put $G_P=F_P/N_P$.
\end{definition}

The free group and quotient group are imported constructions. The normal closure is local: its raw generators are the $d_r$, the identity, products, inverses, and conjugates, followed by propositional truncation.

\formal{GroupPresentation.PresentedGroup.NormalClosure.least; GroupPresentation.PresentedGroup.relationDifferenceInNormalClosure}
\begin{lemma}
$N_P$ is contained in every normal subgroup of $F_P$ containing all $d_r$.
\end{lemma}
\begin{proof}
Eliminate a truncated raw normal-closure witness. A relation generator is in the proposed subgroup by hypothesis; identity, product, and inverse are handled by the subgroup laws; and a conjugate is handled by normality. Hence the constructed subgroup is least. The particular inclusion of $d_r$ is the raw relation constructor.
\end{proof}

\section{The canonical homomorphism}

\formal{GroupPresentation.FundamentalGroup.BG-π₁; GroupPresentation.FundamentalGroup.relationLoopπ₁-null}
Let $\pi_1(B_P)$ denote the fundamental group at the designated base point. The relation loop has trivial class in this group, obtained by taking the set-truncation class of its nullhomotopy.

\formal{GroupPresentation.PresentationHom.freeGroupToBG-π₁; GroupPresentation.PresentationHom.wordToFreeGroup-π₁; GroupPresentation.PresentationHom.relationDifferenceToBG-π₁-null; GroupPresentation.PresentationHom.presentationNormalSubgroupInKernel; GroupPresentation.PresentationHom.canonicalPresentationHom}
\begin{proposition}
There is a canonical group homomorphism $c_P:G_P\to\pi_1(B_P)$.
\end{proposition}
\begin{proof}
Send $g$ to the class of the loop induced by $\ell_g$. The upstream free-group recursor extends this to $F_P\to\pi_1(B_P)$. Induction on a signed word identifies its image with the corresponding concatenated loop class. Hence $d_r$ maps to the class of $L_0(r)L_1(r)^{-1}$, which is trivial by the relation nullhomotopy. By normal-closure leastness, $N_P$ lies in the kernel. The quotient eliminator therefore defines a map on $F_P/N_P$; quotient induction proves multiplicativity. This is $c_P$.
\end{proof}

\section{Surjectivity}

\formal{GroupPresentation.BouquetFundamentalGroup.bouquetLoopHom-surjective; GroupPresentation.Surjectivity.inlPresentation↑-π₁-surjective; GroupPresentation.TruncationNaturality.truncationπ₁-surjective; GroupPresentation.CanonicalSurjectivity.freeGroupToBG-byTruncation≡freeGroupToBG-π₁; GroupPresentation.CanonicalSurjectivity.canonicalPresentationHom-surjective}
\begin{theorem}
For every syntactic presentation $P$, the canonical homomorphism $c_P:G_P\to\pi_1(B_P)$ is surjective.
\end{theorem}
\begin{proof}
The cited bouquet theorem constructs, for every loop class in a bouquet, a free-group word whose image is that class. Its proof uses the library maps \Agda{windingT} and \Agda{loopingT}: after passing through the inverse of the free-groupoid truncation equivalence, the local map \Agda{bouquetLoopHom} agrees with \Agda{loopingT}; the library homotopy \Agda{left-homotopy} makes the recovered word a section. The pointed equivalence of the generator-space section transports this surjectivity to the lifted $X_G$.

The left inclusion $X_G\to Y_P$ is surjective on fundamental groups. Indeed, the cited proof identifies every fibre of $R\times S^1\to R$ with a lifted circle, proves that circle connected, and applies the upstream connected-pushout theorem followed by \Agda{connectedFun→πSurj}. The map $Y_P\to B_P$ induced by the truncation constructor is likewise surjective on fundamental groups: the local proof decodes a truncated loop with \Agda{PathIdTruncIso}, and the inverse law of that isomorphism gives a section.

Composing these three surjections gives a surjection from $F_P$ to $\pi_1(B_P)$. On a free generator, its value is computed by path congruence as the class of the image of $\ell_g$. Free-group homomorphism extensionality therefore identifies this composite with the previously defined free-group map underlying $c_P$. Finally a preimage in $F_P$ is sent to its quotient class in $G_P$; propositional-truncation functoriality transports the equality of images. This is exactly the final proof term.
\end{proof}

\section{Boundary of the result}

The theorem establishes only surjectivity. No declaration proves that the kernel of $c_P$ is trivial; hence neither $G_P\cong\pi_1(B_P)$ nor a Seifert--van Kampen description of this pushout follows from the present development. In particular, although $B_P$ is the natural candidate for a classifying type of the presented group, no pointed equivalence with a classifying space has been formalized. These are mathematical statements requiring additional proofs, not consequences of the notation \Agda{BG_pres}.

\appendix
\section{Formal provenance and reproducibility}

The local results used in the article are proved in \Agda{GroupPresentation.Relations}, \Agda{GroupPresentation.TruncatedRelations}, \Agda{GroupPresentation.Connectedness}, \Agda{GroupPresentation.PresentedGroup}, \Agda{GroupPresentation.PresentationHom}, \Agda{GroupPresentation.GeneratorBouquet}, \Agda{GroupPresentation.BouquetFundamentalGroup}, \Agda{GroupPresentation.Surjectivity}, \Agda{GroupPresentation.TruncationNaturality}, and \Agda{GroupPresentation.CanonicalSurjectivity}. The exact statement-to-source map, including upstream declarations, is maintained in \Agda{alignment/records.jsonl}.

The checked environment is Agda \Agda{2.8.0-3d04bac} with the pinned Cubical checkout \Agda{87af036ebd1c87956cf08fea44e62c7cf76392ca}, on its v0.9-compatible line. The relevant upstream ingredients are the suspension, pushout, truncation, free-group, quotient-group, bouquet fundamental-group, and connected-map interfaces of Cubical Agda. The local proofs cited above were checked with \Agda{agda -i src -i cubical <module>}; all source modules use \AgdaOptions.

The proof of the main theorem uses, rather than re-proves, the upstream bouquet computations \Agda{loopingT}, \Agda{windingT}, and \Agda{left-homotopy}, the connected-map consequence \Agda{connectedFun→πSurj}, and the path/truncation equivalence \Agda{PathIdTruncIso}. Their roles are stated in the proof so that the mathematical argument remains readable independently of this implementation note.


\begin{thebibliography}{9}
\bibitem{hottbook}
The Univalent Foundations Program.
\newblock \emph{Homotopy Type Theory: Univalent Foundations of Mathematics}.
\newblock Institute for Advanced Study, 2013.
\newblock \url{https://homotopytypetheory.org/book/}.

\bibitem{cubical}
Cubical Agda contributors.
\newblock \emph{Cubical Agda standard library}, pinned checkout
\Agda{87af036ebd1c87956cf08fea44e62c7cf76392ca}, v0.9-compatible line.
\newblock \url{https://github.com/agda/cubical}.

\bibitem{agda}
The Agda Team.
\newblock \emph{Agda 2.8.0}.
\newblock \url{https://agda.readthedocs.io/}.
\end{thebibliography}
\end{document}
